% LaTeX template for Artifact Evaluation V20240722, adapted for WiCompass.
% Original template: https://github.com/ctuning/artifact-evaluation/blob/master/docs/template/ae.tex
\documentclass[10pt]{article}
\usepackage[margin=1in]{geometry}
\usepackage{hyperref}
\begin{document}
\special{papersize=8.5in,11in}

\appendix
\section{Artifact Appendix}

\subsection{Abstract}
This artifact supports the evaluation of the MobiCom 2026 paper
``WiCompass: Oracle-driven Data Scaling for mmWave Human Pose Estimation.''
The artifact contains the WiCompass implementation, pre-trained VQ-VAE and
human pose estimation model weights, processed datasets, cached intermediate
results, notebooks for regenerating the paper figures, and scripts for optional
end-to-end re-execution.  The fastest review path uses the supplied model zoo
and experiment logs to regenerate the reported figures.  A deeper review path
can re-run model evaluation, token encoding, KNN coverage, token-space sampling,
pose decoding, simulation-scaling training, real-world training, and pilot-study
training, subject to GPU time availability.

\subsection{Artifact check-list (meta-information)}
{\small
\begin{itemize}
\item {\bf Algorithm:} VQ-VAE pose tokenization; KNN latent-space coverage; prioritized pose sampling (PPS); mmWave human pose estimation with Point Transformer.
\item {\bf Program:} Python package \texttt{wicompass}; experiment scripts and Jupyter notebooks under \texttt{experiments/}; pose-estimation package under \texttt{src/pose\_estimation/}.
\item {\bf Compilation:} Python package installation plus optional PointNet++ CUDA extension compilation.
\item {\bf Transformations:} Human-pose preprocessing; token encoding; KNN index construction; token sampling; token-to-pose decoding; RF simulation traces supplied as datasets.
\item {\bf Binary:} No required binary-only executable.  CUDA extensions are built locally for PointNet++.
\item {\bf Model:} VQ-VAE pose encoder/decoder; Point Transformer based mmWave pose estimation models; released checkpoints in \texttt{model\_zoo/} and \texttt{logs/}.
\item {\bf Data set:} Processed AMASS, MMFi, mmBody, WiCompass real-world data, RF-Genesis simulation datasets, encoded tokens, KNN coverage results, sampled tokens and poses.
\item {\bf Run-time environment:} Linux; Python 3.11; PyTorch 2.x; CUDA 12.x; rclone for Google Drive download; Jupyter for figure notebooks.
\item {\bf Hardware:} NVIDIA CUDA GPU recommended.  At least 16 GB RAM is required; 32 GB RAM is recommended.  Multi-GPU execution is supported for the largest KNN jobs.
\item {\bf Run-time state:} Pre-trained weights and cached logs are provided for short review runs.  Full retraining starts from the supplied datasets.
\item {\bf Execution:} Command-line Python scripts and Jupyter notebooks.
\item {\bf Metrics:} VQ-VAE reconstruction and token metrics; KNN coverage; MPJPE and pose-estimation accuracy metrics; training curves and data-scaling trends.
\item {\bf Output:} Recreated paper figures, JSON result summaries, HDF5 token files, KNN coverage directories, model checkpoints, and notebook outputs.
\item {\bf Experiments:} VQ-VAE microbenchmarks; data coverage; simulator-based data scaling; real-world validation; pilot studies on model size, leave-one-out generalization, and data efficiency.
\item {\bf How much disk space required (approximately)?:} Minimal figure-reproduction path: about 36 GB.  Full workspace: about 83 GB extracted, plus archive/download overhead.
\item {\bf How much time is needed to prepare workflow (approximately)?:} 1--3 hours, dominated by downloading the workspace and compiling CUDA extensions.
\item {\bf How much time is needed to complete experiments (approximately)?:} Figure reproduction from cached logs: minutes to tens of minutes.  Recomputing all KNN coverage needs about 12 GPU-hours.  Full pilot-study retraining needs about 400 GPU-hours on an RTX 3090 class GPU.
\item {\bf Publicly available?:} TODO: add final review URL or public repository URL.
\item {\bf Code licenses (if publicly available)?:} MIT license for the WiCompass repository; third-party submodules retain their own licenses.
\item {\bf Data licenses (if publicly available)?:} TODO: confirm redistribution terms for processed AMASS, MMFi, mmBody, real-world data, and generated simulation data.
\item {\bf Workflow automation framework used?:} Python scripts, shell commands, rclone, Jupyter notebooks, and TensorBoard logs; no external workflow engine is required.
\item {\bf Archived (provide DOI)?:} TODO: add Zenodo/ACM DL/other archival DOI if requesting Artifacts Available.
\end{itemize}
}

\subsection{Description}

\subsubsection{How to access}
TODO: provide the final artifact package link for reviewers.  The working code
repository is expected to contain the WiCompass source tree with the following
main components:
\begin{itemize}
\item \texttt{src/wicompass/}: VQ-VAE, latent-space coverage, sampling, evaluation, and visualization code.
\item \texttt{src/pose\_estimation/}: mmWave pose-estimation datasets, model, training, and evaluation code.
\item \texttt{experiments/}: scripts and notebooks that regenerate the paper figures and optional training results.
\item \texttt{tools/}: workspace download, extraction, visualization, and setup utilities.
\end{itemize}

The data, logs, and model weights are distributed as a WiCompass workspace.  The
workspace can be downloaded using rclone from the shared Google Drive folder:
\begin{center}
\url{https://drive.google.com/drive/folders/1GDgcJ-6fq4TW-AmuZPPa-i_UvshAEg79?usp=sharing}
\end{center}
The command \texttt{python tools/download\_workspace.py --minimal} downloads the
minimal artifact components for reproducing figures from cached logs.  The
command \texttt{python tools/download\_workspace.py} downloads all components.
After extraction, \texttt{python tools/setup\_workspace.py --workspace <path>}
creates repository-local symbolic links.

\subsubsection{Hardware dependencies}
A CUDA-capable NVIDIA GPU is recommended for model evaluation, KNN coverage, and
training.  The authors tested the code path with CUDA 12.x and PyTorch 2.x.  The
largest experiments benefit from multiple GPUs; however, the cached-log path for
figure regeneration can be inspected on a smaller GPU machine or CPU-only
machine when notebooks only parse existing JSON/HDF5 results.

\subsubsection{Software dependencies}
The artifact assumes Linux, Conda, Python 3.11, CUDA 12.x, PyTorch 2.x, and
rclone.  Core Python dependencies include \texttt{torch}, \texttt{torchvision},
\texttt{numpy}, \texttt{scipy}, \texttt{faiss-gpu-cu12}, \texttt{h5py},
\texttt{pandas}, \texttt{PyYAML}, \texttt{scikit-learn}, \texttt{matplotlib},
\texttt{seaborn}, \texttt{tensorboard}, \texttt{einops}, \texttt{ray},
\texttt{loguru}, \texttt{omegaconf}, \texttt{smplx}, \texttt{transforms3d},
\texttt{human\_body\_prior}, and \texttt{chumpy}.  The \texttt{chumpy} package
should be installed before \texttt{human\_body\_prior}.  PointNet++ CUDA
extensions are built from the repository submodule.

\subsubsection{Data sets}
The workspace contains the following extracted components:
\begin{itemize}
\item \texttt{model\_zoo}: pre-trained model weights, about 1.6 GB.
\item \texttt{wicompass\_logs}: training logs, encoded tokens, KNN results, and cached figure inputs, about 34 GB.
\item \texttt{AMASS\_preproc}: processed AMASS pose data, about 3.2 GB.
\item \texttt{mmfi}: MMFi dataset subset used by the artifact, about 1.6 GB.
\item \texttt{mmbody}: mmBody dataset subset used by the artifact, about 34 GB.
\item \texttt{real\_world}: WiCompass real-world collected data, about 431 MB.
\item \texttt{simulation\_datasets}: RF-Genesis generated mmWave simulation datasets, about 8.2 GB.
\end{itemize}

\subsubsection{Models}
The main supplied model artifacts are the VQ-VAE checkpoint used for pose
encoding and decoding, pose-estimation checkpoints used in the simulator and
real-world evaluations, and pilot-study checkpoints used for model-size,
leave-one-out, and data-efficiency comparisons.  The VQ-VAE checkpoint used in
coverage analysis is stored at
\texttt{logs/vqvae/vqvae\_tokennum16\_tokenclass64/best\_model.pth}.

\subsection{Installation}
Create and activate the Conda environment:
\begin{verbatim}
conda create -n wicompass python=3.11 -y
conda activate wicompass
\end{verbatim}
Install Python dependencies and the WiCompass package:
\begin{verbatim}
pip install --no-build-isolation chumpy
pip install -r requirements.txt
pip install -e .
\end{verbatim}
Build the PointNet++ CUDA extension:
\begin{verbatim}
git submodule update --init --recursive
cd submodules/Pointnet2_PyTorch/pointnet2_ops_lib
export CC=/usr/bin/gcc
export CXX=/usr/bin/g++
pip install -e . --no-build-isolation
\end{verbatim}
Download and link the workspace:
\begin{verbatim}
python tools/download_workspace.py --workspace ~/wicompass_workspace --minimal
python tools/setup_workspace.py --workspace ~/wicompass_workspace
\end{verbatim}
For full retraining or full pipeline inspection, omit \texttt{--minimal} to
retrieve all dataset components.

\subsection{Experiment workflow}
The recommended review workflow has three levels.

\paragraph{Level 1: Fast figure reproduction from cached logs.}
Download the minimal workspace and run the notebooks that parse released JSON,
HDF5, and log files:
\begin{itemize}
\item Figure 10: \texttt{experiments/vqvae\_microbenchmark/vqvae\_performance.ipynb}.
\item Figures 7 and 8: \texttt{experiments/knn\_coverage/intra\_knn\_distribution.ipynb} and \texttt{experiments/knn\_coverage/knn\_coverage.ipynb}.
\item Figure 6: \texttt{experiments/simulation\_scaling/simulation\_results.ipynb}.
\item Figure 9: \texttt{experiments/real\_world\_scaling/real\_world\_paper\_figure.ipynb}.
\item Figure 1 and pilot studies: notebooks under \texttt{experiments/pilot\_study/}.
\end{itemize}

\paragraph{Level 2: Recompute selected intermediate results.}
Reviewers can evaluate supplied checkpoints and regenerate intermediate files:
\begin{verbatim}
cd experiments/vqvae_microbenchmark
python recover_training_results.py
python batch_evaluate_models.py
\end{verbatim}
Encode datasets into tokens:
\begin{verbatim}
python src/wicompass/evaluation/encode_datasets_into_tokens.py \
  --dataset MMFi --dataset-type mmfi
python src/wicompass/evaluation/encode_datasets_into_tokens.py \
  --dataset MMBody --dataset-type mmbody
\end{verbatim}
Recompute KNN coverage for selected settings:
\begin{verbatim}
python src/wicompass/knn_coverage/knn_coverage.py \
  --dataset-A logs/wicompass/encoded_tokens/BMLmovi-BMLrub-CMU-GRAB-KIT-MOYO-MoSh-PosePrior-WEIZMANN_tokens.h5 \
  --dataset-B logs/wicompass/encoded_tokens/MMFi_tokens.h5 \
  --metric cosine --k 2 --multi-gpu --sample-ratio 1.0 \
  --output logs/wicompass/knn_coverage/mmfi/k2
\end{verbatim}

\paragraph{Level 3: Re-run training and data-scaling experiments.}
With the full workspace, reviewers can re-run simulator-scaling training:
\begin{verbatim}
python experiments/simulation_scaling/train_simulation_batch.py --list
python experiments/simulation_scaling/train_simulation_batch.py
python experiments/simulation_scaling/train_simulation_batch_different_size.py
\end{verbatim}
They can also re-run real-world training and pilot-study training scripts under
\texttt{experiments/real\_world\_scaling/} and \texttt{experiments/pilot\_study/}.
Full pilot-study retraining is expected to require about 400 GPU-hours on an
RTX 3090 class GPU.

\subsection{Evaluation and expected results}
For the fast review path, the expected result is a regenerated set of plots that
matches the trends and main claims in the paper: VQ-VAE model configurations
produce the reported reconstruction/tokenization tradeoffs, WiCompass improves
latent-space coverage over baseline mmWave datasets, WiCompass-selected data
improves simulator-based pose-estimation scaling, the real-world validation
follows the paper's reported trend, and pilot studies reproduce the motivating
model-size, leave-one-out, and data-efficiency observations.  Exact pixel-level
or floating-point equality is not expected because plotting libraries, GPU
kernels, and stochastic training can vary across systems; differences should not
change the paper's conclusions.

For intermediate recomputation, expected outputs include \texttt{training\_result.json},
\texttt{testing\_result.json}, HDF5 token files in \texttt{logs/wicompass/encoded\_tokens/},
KNN coverage outputs under \texttt{logs/wicompass/knn\_coverage/}, sampled token
arrays under \texttt{logs/wicompass/sampled\_tokens/}, and sampled pose outputs
under \texttt{logs/wicompass/sampled\_poses/}.

\subsection{Experiment customization}
The artifact can be customized by changing VQ-VAE configuration files under
\texttt{src/wicompass/configs/}, selecting different datasets in
\texttt{src/wicompass/configs/datasets.json}, varying KNN parameters such as
\texttt{k}, metric, and sample ratio in \texttt{src/wicompass/knn\_coverage/},
and adjusting sampling budgets, seeds, and coverage parameters in
\texttt{src/wicompass/token\_space\_sampling/}.  Pose-estimation training can be
customized through YAML files under \texttt{src/pose\_estimation/configs/} and
\texttt{experiments/real\_world\_scaling/configs/}.

\subsection{Notes}
The artifact has a short path and a full path.  The short path is intended for
artifact reviewers who need to validate documentation, executability, and the
main reported trends within limited review time.  The full path is intended for
reviewers with access to substantial GPU time and storage.  RF-Genesis is used
to produce simulated mmWave data, but the generated simulation datasets are
included so reviewers do not need to install RF-Genesis unless they want to audit
that external simulation stage.

TODO: verify the target venue's anonymity policy.  If the AE process is double
anonymous, replace author-identifying URLs and repository names with anonymous
review links until de-anonymization is allowed.

\subsection{Methodology}
This appendix is organized according to the cTuning Artifact Evaluation template
and ACM artifact badging terminology.  The submission primarily targets
Artifacts Evaluated -- Functional, because the artifact is documented,
consistent with the paper, complete to the extent permitted by dataset licenses,
and exercisable through released scripts, notebooks, model weights, and cached
results.  If the final artifact is deposited in a permanent archival repository
with a DOI or equivalent unique identifier, it can also request Artifacts
Available.  Results Validated should only be requested after an independent
non-author team has reproduced or replicated the main results.

\begin{itemize}
\item \url{https://www.acm.org/publications/policies/artifact-review-and-badging-current}
\item \url{https://cTuning.org/ae}
\end{itemize}

\end{document}
